\documentclass[10pt,a4paper]{amsart}
\usepackage{amsmath,amssymb,amsthm,graphicx}
% Retained minimal preamble for the Stacks Project chapter flat.tex (More on Flatness).
% Verbatim excerpt of the chapter preamble: Xy-pic setup (native preamble lines 14-18) and the native
% theorem declarations with the mathematical macros (native preamble lines 154-217). Copied mechanically;
% nothing added, removed or redefined. The wrapper class and the amsmath/amssymb/amsthm/graphicx base
% packages are supplied by the assembler, so the class/externaldocument part of the native preamble is omitted.
% For commutative diagrams we use Xy-pic
\usepackage[all]{xy}

% We use 2cell for 2-commutative diagrams.
\xyoption{2cell}
% Theorem environments.
%
\theoremstyle{plain}
\newtheorem{theorem}[subsection]{Theorem}
\newtheorem{proposition}[subsection]{Proposition}
\newtheorem{lemma}[subsection]{Lemma}

\theoremstyle{definition}
\newtheorem{definition}[subsection]{Definition}
\newtheorem{example}[subsection]{Example}
\newtheorem{exercise}[subsection]{Exercise}
\newtheorem{situation}[subsection]{Situation}

\theoremstyle{remark}
\newtheorem{remark}[subsection]{Remark}
\newtheorem{remarks}[subsection]{Remarks}

\numberwithin{equation}{subsection}

% Macros
%
\def\lim{\mathop{\mathrm{lim}}\nolimits}
\def\colim{\mathop{\mathrm{colim}}\nolimits}
\def\Spec{\mathop{\mathrm{Spec}}}
\def\Hom{\mathop{\mathrm{Hom}}\nolimits}
\def\Ext{\mathop{\mathrm{Ext}}\nolimits}
\def\SheafHom{\mathop{\mathcal{H}\!\mathit{om}}\nolimits}
\def\SheafExt{\mathop{\mathcal{E}\!\mathit{xt}}\nolimits}
\def\Sch{\mathit{Sch}}
\def\Mor{\mathop{\mathrm{Mor}}\nolimits}
\def\Ob{\mathop{\mathrm{Ob}}\nolimits}
\def\Sh{\mathop{\mathit{Sh}}\nolimits}
\def\NL{\mathop{N\!L}\nolimits}
\def\CH{\mathop{\mathrm{CH}}\nolimits}
\def\proetale{{pro\text{-}\acute{e}tale}}
\def\etale{{\acute{e}tale}}
\def\QCoh{\mathit{QCoh}}
\def\Ker{\mathop{\mathrm{Ker}}}
\def\Im{\mathop{\mathrm{Im}}}
\def\Coker{\mathop{\mathrm{Coker}}}
\def\Coim{\mathop{\mathrm{Coim}}}

% Boxtimes
%
\DeclareMathSymbol{\boxtimes}{\mathbin}{AMSa}{"02}

%
% Macros for moduli stacks/spaces
%
\def\QCohstack{\mathcal{QC}\!\mathit{oh}}
\def\Cohstack{\mathcal{C}\!\mathit{oh}}
\def\Spacesstack{\mathcal{S}\!\mathit{paces}}
\def\Quotfunctor{\mathrm{Quot}}
\def\Hilbfunctor{\mathrm{Hilb}}
\def\Curvesstack{\mathcal{C}\!\mathit{urves}}
\def\Polarizedstack{\mathcal{P}\!\mathit{olarized}}
\def\Complexesstack{\mathcal{C}\!\mathit{omplexes}}
% \Pic is the operator that assigns to X its picard group, usage \Pic(X)
% \Picardstack_{X/B} denotes the Picard stack of X over B
% \Picardfunctor_{X/B} denotes the Picard functor of X over B
\def\Pic{\mathop{\mathrm{Pic}}\nolimits}
\def\Picardstack{\mathcal{P}\!\mathit{ic}}
\def\Picardfunctor{\mathrm{Pic}}
\def\Deformationcategory{\mathcal{D}\!\mathit{ef}}

% --- appended 2026-10-12 (Chief): the project own definitions of the reference/slogan/history
% display environments, copied mechanically from the shared Stacks preamble lines 8-12
% (\usepackage{verbatim} plus the three \newenvironment lines, which the chapter preamble keeps
% verbatim from the project preamble). Nothing added, removed or redefined: these environments
% carry no mathematical content and the project itself typesets them as comments. Without these
% three lines a counted statement whose own text uses them fails with Environment slogan undefined.
\usepackage{verbatim}
\newenvironment{reference}{\comment}{\endcomment}
\newenvironment{slogan}{\comment}{\endcomment}
\newenvironment{history}{\comment}{\endcomment}
\usepackage{cleveref}
\crefname{originallemma}{original lemma}{original lemmas}
\Crefname{originallemma}{Original lemma}{Original lemmas}
\makeatletter
\newlabel{lemma-finite-type-flat-along-fibre-free}{{05L0}{source}{}{sourcepointer-1}{}}
\newlabel{lemma-finite-type-flat-along-fibre-free@cref}{{[originallemma][0][]05L0}{[1][0][]source}{}{}{}}
\makeatother
\crefname{originallemma}{original lemma}{original lemmas}
\Crefname{originallemma}{Original lemma}{Original lemmas}
\makeatletter
\newlabel{lemma-finite-type-flat-at-point-free-variant}{{05KZ}{source}{}{sourcepointer-2}{}}
\newlabel{lemma-finite-type-flat-at-point-free-variant@cref}{{[originallemma][0][]05KZ}{[1][0][]source}{}{}{}}
\makeatother
\crefname{originallemma}{original lemma}{original lemmas}
\Crefname{originallemma}{Original lemma}{Original lemmas}
\makeatletter
\newlabel{more-morphisms-lemma-elementary-etale-neighbourhoods}{{057B}{source}{}{sourcepointer-3}{}}
\newlabel{more-morphisms-lemma-elementary-etale-neighbourhoods@cref}{{[originallemma][0][]057B}{[1][0][]source}{}{}{}}
\makeatother
\title{The flat along the fibre freeness statement without the finite presentation hypothesis}
\author{The Stacks Project Authors}
\date{stacks-project:flat.tex:local-snapshot-20261011; GNU FDL 1.2 or later}
\begin{document}
\maketitle
\noindent\textit{Editorial note.} All delivered mathematical text is the author's own, copied from the pinned Stacks Project chapter file flat.tex; no statement or proof step was rewritten, reordered or removed.\par\medskip
\noindent\textit{Editorial note.} Preamble provenance: the chapter's own retained minimal preamble (sources/stacks-flat/retained-preamble.tex, lines 1-85) is a verbatim excerpt of the native Xy-pic setup and of the native theorem declarations and mathematical macros, followed by the project's own reference/slogan/history display-environment definitions copied mechanically from the shared preamble; the wrapper class and the amsmath/amssymb/amsthm/graphicx base packages are supplied by the assembler, and extra\_packages is left empty because everything needed is already in the retained preamble. No macro was invented or redefined.\par\medskip
\noindent\textit{Editorial note.} Same-chapter pointers are located at their \textbackslash{}label lines inside the pinned chapter file; cross-chapter pointers are located in the project's own tag list staged in the snapshot (sources/stacks-chow/tags.txt), at the verbatim lines carrying their canonical tags.\par\medskip
\noindent\textit{Allowed standard prerequisites (editorial).}\begin{itemize}
\item The pointwise free version of the flat along the fibre lemma in its finite type form. Same chapter flat.tex, label lemma-finite-type-flat-at-point-free-variant, invoked at proof line 3656. Produces, for each point of the fibre, an elementary etale neighbourhood and an affine etale neighbourhood whose global sections of the pulled back module are free over the base, with residue fields preserved; assembling finitely many of them is what the counted proof does.
\item Existence of an elementary etale neighbourhood dominating finitely many given ones. More on Morphisms, label more-morphisms-lemma-elementary-etale-neighbourhoods, invoked at proof line 3676. Supplies the single dominating elementary etale neighbourhood of the base which replaces the finitely many pointwise ones.
\item The finite presentation version of the same statement, used as the model for this one. Same chapter flat.tex, label lemma-finite-type-flat-along-fibre-free, invoked at proof line 3653 and claimed separately as S14-algebra-flat-along-fibre-free. The author notes in the first sentence of the proof that the only difference from that lemma is the absence of the finite presentation hypothesis on the morphism.
\end{itemize}\medskip
\ifdefined\hypertarget\hypertarget{sourcepointer-1}{}\fi
\noindent\textit{Source reference.} Original lemma 05L0: sources/stacks-flat/pinned-original-source.tex, source lines 3557--3557. Navigation pointer for \textbackslash{}ref\{lemma-finite-type-flat-along-fibre-free\} used inside the counted span (line 3653); the argument that consumes it is retained in the span.\par\medskip
\ifdefined\hypertarget\hypertarget{sourcepointer-2}{}\fi
\noindent\textit{Source reference.} Original lemma 05KZ: sources/stacks-flat/pinned-original-source.tex, source lines 3501--3501. Navigation pointer for \textbackslash{}ref\{lemma-finite-type-flat-at-point-free-variant\} used inside the counted span (line 3656); the argument that consumes it is retained in the span.\par\medskip
\ifdefined\hypertarget\hypertarget{sourcepointer-3}{}\fi
\noindent\textit{Source reference.} Original lemma 057B: sources/stacks-chow/tags.txt, source lines 6739--6739. Navigation pointer for \textbackslash{}ref\{more-morphisms-lemma-elementary-etale-neighbourhoods\} used inside the counted span (line 3676); the argument that consumes it is retained in the span.\par\medskip
\begin{lemma}
\label{lemma-finite-type-flat-along-fibre-free-variant}
Let $f : X \to S$ be a morphism of schemes.
Let $\mathcal{F}$ be a quasi-coherent sheaf on $X$.
Let $s \in S$.
Assume that
\begin{enumerate}
\item $f$ is of finite type,
\item $\mathcal{F}$ is of finite type, and
\item $\mathcal{F}$ is flat over $S$ at all points of $X_s$.
\end{enumerate}
Then there exists an elementary \'etale neighbourhood $(S', s') \to (S, s)$
and a commutative diagram of schemes
$$
\xymatrix{
X \ar[d] & X' \ar[l]^g \ar[d] \\
S & \Spec(\mathcal{O}_{S', s'}) \ar[l]
}
$$
such that $X' \to X \times_S \Spec(\mathcal{O}_{S', s'})$
is \'etale, $X_s = g((X')_{s'})$, the scheme $X'$ is affine,
and such that $\Gamma(X', g^*\mathcal{F})$ is a free
$\mathcal{O}_{S', s'}$-module.
\end{lemma}

\begin{proof}
(The only difference with
Lemma \ref{lemma-finite-type-flat-along-fibre-free}
is that we do not assume $f$ is of finite presentation.)
For every point $x \in X_s$ we can use
Lemma \ref{lemma-finite-type-flat-at-point-free-variant}
to find an elementary \'etale neighbourhood $(S_x , s_x) \to (S, s)$
and a commutative diagram
$$
\xymatrix{
(X, x) \ar[d] & (Y_x, y_x) \ar[l]^{g_x} \ar[d] \\
(S, s) & (\Spec(\mathcal{O}_{S_x, s_x}), s_x) \ar[l]
}
$$
such that $Y_x \to X \times_S \Spec(\mathcal{O}_{S_x, s_x})$
is \'etale, $\kappa(x) = \kappa(y_x)$, the scheme $Y_x$ is affine, and
such that $\Gamma(Y_x, g_x^*\mathcal{F})$ is a free
$\mathcal{O}_{S_x, s_x}$-module. In particular
$g_x((Y_x)_{s_x})$ is an open neighbourhood of $x$ in $X_s$.
Because $X_s$ is quasi-compact we can find a finite number of points
$x_1, \ldots, x_n \in X_s$ such that $X_s$ is the union of
the $g_{x_i}((Y_{x_i})_{s_{x_i}})$. Choose an elementary \'etale neighbourhood
$(S' , s') \to (S, s)$ which dominates each of the neighbourhoods
$(S_{x_i}, s_{x_i})$, see
More on Morphisms,
Lemma \ref{more-morphisms-lemma-elementary-etale-neighbourhoods}.
Set
$$
X' = \coprod Y_{x_i} \times_{\Spec(\mathcal{O}_{S_{x_i}, s_{x_i}})}
\Spec(\mathcal{O}_{S', s'})
$$
and endow it with the obvious morphism $g : X' \to X$.
By construction $X_s = g(X'_{s'})$ and
$$
\Gamma(X', g^*\mathcal{F})
=
\bigoplus \Gamma(Y_{x_i}, g_{x_i}^*\mathcal{F})
\otimes_{\mathcal{O}_{S_{x_i}, s_{x_i}}}
\mathcal{O}_{S', s'}.
$$
This is a free $\mathcal{O}_{S', s'}$-module as a direct sum
of base changes of free modules.
\end{proof}

\end{document}
